\section{Keep one working source through the construction}
Work in one copy of \path{packages/029-binary-ppm/starter/main.asm}. It supplies the caller,
object arenas, checkpoint labels and frame/cleanup paths. Each named BEGIN/END region
is initially unfinished. Author only the requested body, keeping those markers and the
labels needed by other regions. The first caller mode selects parsing of the worked
memory file; it returns development status four before any implementation is present.

Development status is outside the completed interface's statuses zero through three.
It does not certify the bytes already produced by earlier regions. During construction,
a prefix or even a complete image can be present before the unfinished return becomes
a successful one. Observe those bytes at their checkpoints, then inspect the actual
returned status separately. The caller's \texttt{count} field retains all 64 returned
bits, including both dimensions of a decoder result; its unit depends on the selected
interface.

Use \texttt{python3 prepare.py} in the construction directory with an exact recipe ID,
\texttt{--source} naming your current working file and \texttt{--output} naming a new
observation file. Preparation changes data arenas and caller controls only. It retains
authored instructions and refuses to overwrite an existing output. For example,
\texttt{s02-1} prepares the first count check. In the supplied caller, stop at the
recipe's \texttt{callerCheckpoint}. A direct-call \texttt{function\_return} means the
entry-specific caller return label when using \texttt{main}.

File recipes additionally describe fixture contents, short-operation limits, error
positions and descriptor selection. Preparing assembly does not install an in-memory
file environment in an ordinary IDE run. The guided verifier supplies that environment;
for a normal successful Run or Debug, the corresponding input file must exist in the
configured working directory. Automated capture provisioning remains subsequent work.
An injected short-write or error recipe is a grading-environment observation, not a
claim about a host file service producing that same result on demand.

Guided checks inspect the illustrated intermediate state. Completed free-writing checks
accept equivalent registers, frames and algorithms that satisfy the published interface.
A private implementation is not an instruction template that students must match.

\section{Read the supplied frame and cleanup paths}
Nested decoder, loader and saver routines use a 128-byte frame. The supplied saves and
restores retain every bit of \texttt{s0} through \texttt{s7}, \texttt{fp} and
\texttt{ra}. Reserve bytes 80--95 for a private metadata descriptor and byte 96 for the
loader's separate EOF probe. The public metadata object and the private descriptor are
different storage, even when their contents will agree on success.
\tagpdfsetup{table/header-rows={1}}
\begin{center}
\begin{tabular}{rl}
Frame offset & Supplied use \\
\hline
0--56 & Eight saved registers, eight bytes each \\
64 & Saved frame pointer \\
72 & Saved return address \\
80--95 & Private metadata descriptor \\
96 & Private one-byte EOF probe \\
\end{tabular}
\end{center}

The guided file routines use \texttt{s3} for their live descriptor. The supplied frame
initializes it to minus one before an open is attempted. Keep that marker until a
successful open replaces it. The development exit closes a nonnegative descriptor,
chooses status four and zero count, then restores the frame. Completed file errors
have their own close-and-error paths. Completing the final load/save body reaches the
supplied restore rather than falling through to the development exit.
